\documentclass[11pt]{article}
\usepackage[T1]{fontenc}
\usepackage{lmodern}
\usepackage{amsmath,amssymb,amsthm}
\usepackage{enumitem}
\usepackage{graphicx}
\usepackage[margin=0.92in]{geometry}
\usepackage[colorlinks=true,linkcolor=blue,citecolor=blue,urlcolor=blue]{hyperref}

\newtheorem{theorem}{Theorem}
\newtheorem{lemma}{Lemma}
\newtheorem{remark}{Remark}
\newcommand{\dv}{\,dv_g}
\newcommand{\R}{\mathbb R}
\newcommand{\calE}{\mathcal E}
\newcommand{\calB}{\mathcal B}
\newcommand{\Ltwo}{L^2(M,v_g)}

\title{Joint Smooth Approximation of Compact Weighted Traces\\
and Convergence of All Variational Eigenvalues}
\author{Run \texttt{fa\_banach\_001}, agent \texttt{agent\_lane\_14}}
\date{19 July 2026}

\begin{document}
\maketitle

\section*{Status and scope}

Candidate full result, likely valid, for the uniformly positive setting
explicitly invoked immediately before Remark~1.3 of Vinokurov's
arXiv:2601.17896v3~\cite{Vinokurov}.  The result removes the continuity and
upper-boundedness assumptions from the earlier partial argument: the energy
density may be rough and unbounded.  It proves convergence, not merely the
requested upper bound, for every eigenvalue at once.

The theorem assumes \(\alpha\geq c>0\), exactly as in the sentence preceding
the open question.  It does not treat a different, broader interpretation in
which \(\alpha\) may vanish.  That distinction matters: the source notes that
for degenerate weights the completed space \(H^1(\alpha)\) can have nonunique
differentials and even nonzero pure-differential elements.

\section*{Source question}

Vinokurov asks in Remark~1.3, PDF page~4, whether the eigenvalues associated
with an arbitrary measure inducing a compact identity
\(H^1(M,\alpha)\to L^2(M,\mu)\) can be approximated, or at least bounded from
above, by eigenvalues associated with smooth measures~\cite{Vinokurov}.  The
source passage, including the preceding hypothesis \(\alpha\geq c>0\), is
reproduced below.

\begin{center}
\includegraphics[width=\linewidth]{figures/open_problem_crop.png}
\end{center}

\section*{Definitions}

Let \((M,g)\) be a closed connected Riemannian manifold.  Let
\(\alpha\in L^1(M)\) satisfy \(\alpha\geq c>0\) almost everywhere, and let
\(\mu\neq0\) be a finite positive Radon measure.  Write
\[
 a_\alpha(u,v)=\int_M\alpha\langle du,dv\rangle_g\dv,
 \qquad b_\mu(u,v)=\int_M uv\,d\mu.
\]
The space \(H_\alpha=H^1(M,1,\alpha)\) is the completion of smooth functions
for
\[
 \|u\|_{H_\alpha}^2=\|u\|_{\Ltwo}^2+a_\alpha(u,u).
\]
Uniform positivity makes the function and weak differential components
unambiguous and gives a continuous injection \(H_\alpha\hookrightarrow
H^1(M)\).  We assume that the trace
\[
 J:H_\alpha\longrightarrow L^2(M,\mu)
\]
extends compactly from smooth functions.

The variational eigenvalues are
\[
 \lambda_k(\alpha,\mu)=
 \inf_{\substack{F\subset C^\infty(M)\\ \dim_{L^2(\mu)}F=k+1}}
 \ \sup_{0\neq u\in F}\frac{a_\alpha(u,u)}{b_\mu(u,u)},
 \qquad k=0,1,\ldots .
\]
Compactness of \(J\), together with the Poincar\'e inequality, gives the usual
discrete generalized spectrum, with \(\lambda_0=0\).

For \(2<p<\infty\), put \(q=p/(p-2)\).  When
\(\alpha\in L^q\), Vinokurov's normalization is
\[
 \overline\lambda_{k,p}(\alpha,\mu)
 =\lambda_k(\alpha,\mu)\frac{\mu(M)}{\|\alpha\|_{L^q}}.
\]
For \(p=2\), the source fixes \(\alpha\equiv1\).

\section*{Proof intuition}

Smoothing \(\mu\) alone has the wrong one-sided min--max behavior when
\(\alpha\) is rough.  The key is to smooth the measures \(\alpha v_g\) and
\(\mu\) by the same small diffeomorphisms.

Compactness of the limiting trace is equivalent to an Ehrling-type estimate
\[
 b_\mu(v,v)\leq\varepsilon a_\alpha(v,v)
       +C_\varepsilon\|v\|_{\Ltwo}^2.
\]
Apply it to \(v=u\circ\Phi_z\), then average in the smoothing parameter.
The averaged mass is exactly the smoothed mass form.  The averaged energy is,
up to \(1+O(|z|)\), the energy of the simultaneously smoothed scalar density.
Thus the entire smooth family inherits one uniform compactness estimate.

Meanwhile, the smoothed weights converge strongly in \(L^1\), and their
energy forms converge in the Mosco sense.  The uniform compactness estimate
then says that strong \(L^2(v_g)\) convergence plus bounded moving energy
forces convergence in the moving \(L^2(\mu_n)\) norms.  Finite-dimensional
recovery spaces give the eigenvalue limsup; limits of the first \(k+1\)
smooth eigenfunctions give the liminf.

\section*{Main result}

\begin{theorem}[joint smoothing of every compact weighted trace]
\label{thm:main}
Let \((M,g)\) be closed and connected, let
\(\alpha\in L^1(M)\) satisfy \(\alpha\geq c>0\) almost everywhere, and let
\(\mu\neq0\) be a finite positive Radon measure such that
\[
 H^1(M,1,\alpha)\longrightarrow L^2(M,\mu)
\]
is compact.  There are \(\alpha_n,\rho_n\in C^\infty(M)\), both strictly
positive, such that, with \(\mu_n=\rho_n v_g\),
\begin{align*}
 &\alpha_n\longrightarrow\alpha &&\text{in }L^1(M),\\
 &\mu_n\stackrel{*}{\rightharpoonup}\mu,
   \qquad \mu_n(M)=\mu(M),\\
 &\lambda_k(\alpha_n,\mu_n)\longrightarrow
   \lambda_k(\alpha,\mu) &&\text{for every }k\geq0.
\end{align*}
If \(\alpha\in L^r(M)\) for some \(1\leq r<\infty\), the construction has
\(\alpha_n\to\alpha\) in \(L^r\).  Consequently, for \(2<p<\infty\) and
\(\alpha\in L^{p/(p-2)}\), every normalized eigenvalue also converges.  In
the \(p=2\) case with \(\alpha=1\), one may take \(\alpha_n\equiv1\).
\end{theorem}

In particular, for \(2<p\leq m\), every pair in the uniformly positive
regime of Remark~1.3 satisfies
\[
 \overline\lambda_{k,p}(\alpha,\mu)
 \leq \sup_{\substack{\widetilde\alpha,\widetilde\mu\in C^\infty(M)\\
                 \widetilde\alpha,\widetilde\mu>0}}
       \overline\lambda_{k,p}(\widetilde\alpha,
                              \widetilde\mu v_g),
\]
which is the requested smooth upper bound; the theorem gives the stronger
approximation statement.

\section*{Proof}

\subsection*{1. Compactness gives an Ehrling estimate}

\begin{lemma}\label{lem:ehrling}
For every \(\varepsilon>0\), there is \(C_\varepsilon<\infty\) such that
\begin{equation}\label{eq:ehrling}
 b_\mu(v,v)\leq \varepsilon a_\alpha(v,v)
      +C_\varepsilon\|v\|_{\Ltwo}^2
 \qquad(v\in H_\alpha).
\end{equation}
\end{lemma}

\begin{proof}
If this failed for some \(\varepsilon\), after normalization there would be
\(v_j\in H_\alpha\) with
\[
 b_\mu(v_j,v_j)=1,\qquad
 a_\alpha(v_j,v_j)\leq\varepsilon^{-1},\qquad
 \|v_j\|_{\Ltwo}^2\longrightarrow0.
\]
The sequence is bounded in \(H_\alpha\), and any weakly convergent
subsequence has weak limit zero because its \(L^2(v_g)\) norm tends to zero.
Compactness of \(J\) would then give
\(b_\mu(v_j,v_j)\to0\), a contradiction.
\end{proof}

\subsection*{2. A common intrinsic smoothing family}

Choose finitely many smooth global vector fields
\(X_1,\ldots,X_N\) whose values span \(T_xM\) at every \(x\).  Such a
family is obtained, for example, by embedding \(M\) in Euclidean space and
projecting a basis of the ambient space onto \(TM\).  Let
\[
 \Phi_z=\operatorname{Fl}^{X_N}_{z_N}\circ\cdots\circ
          \operatorname{Fl}^{X_1}_{z_1},
 \qquad z\in\R^N
\]
for small \(z\).  Then \(\Phi_0=\mathrm{id}\), the maps are uniformly
bi-Lipschitz, and
\begin{equation}\label{eq:metric-close}
 (1-C|z|)g\leq \Phi_z^*g\leq(1+C|z|)g.
\end{equation}
Moreover \(z\mapsto\Phi_z(x)\) is a submersion near \(z=0\), uniformly in
\(x\); these are standard consequences of smooth flows and the submersion
theorem~\cite{Lee}.

Let \(\eta_t\) be a nonnegative smooth density on \(\R^N\), supported in
\(|z|<t\), with integral one.  For a finite measure \(\sigma\) on \(M\), set
\begin{equation}\label{eq:Pt}
 P_t\sigma=\int_{\R^N}(\Phi_z)_*\sigma\,\eta_t(z)\,dz.
\end{equation}
The submersion theorem and integration along the parameter fibers give a
smooth kernel \(K_t(x,y)\) such that
\[
 d(P_t\sigma)(y)=\left(\int_MK_t(x,y)\,d\sigma(x)\right)dv_g(y).
\]
Thus \(P_t\sigma\) has a smooth density even when \(\sigma\) is singular.
It preserves total mass and converges weak-star to \(\sigma\).  If
\(d\sigma=f\,dv_g\) with \(f\in L^r\), \(1\leq r<\infty\), strong continuity
of pullback by a smooth family of diffeomorphisms gives
\begin{equation}\label{eq:Lr}
 \frac{dP_t(fv_g)}{dv_g}\longrightarrow f\quad\text{in }L^r(M).
\end{equation}

Define
\[
 \alpha_t\,dv_g=P_t(\alpha v_g),\qquad \nu_t=P_t\mu.
\]
The first density is smooth.  Since \(\alpha\geq c\) and the Jacobians of
\(\Phi_z\) equal \(1+O(t)\) uniformly,
\begin{equation}\label{eq:lower}
 \alpha_t\geq c/2
\end{equation}
for all sufficiently small \(t\).  To make the mass density strictly positive
while preserving mass, choose \(\delta_t\downarrow0\) and put
\begin{equation}\label{eq:positive-mass}
 \mu_t=(1-\delta_t)\nu_t+
 \delta_t\frac{\mu(M)}{v_g(M)}v_g.
\end{equation}

For smooth \(u\), put \(U_zu=u\circ\Phi_z\).  The definitions give
\begin{align}
 b_{\nu_t}(u,u)&=\int\!b_\mu(U_zu,U_zu)\eta_t(z)\,dz,
 \label{eq:mass-average}\\
 a_{\alpha_t}(u,u)&=\int\!\int_M
   \alpha(x)|du(\Phi_z(x))|_g^2\dv(x)\eta_t(z)\,dz.
 \label{eq:energy-scalar}
\end{align}
By \eqref{eq:metric-close}, if
\(A_t(u)=\int a_\alpha(U_zu,U_zu)\eta_t(z)\,dz\), then
\begin{equation}\label{eq:energy-compare}
 (1-Ct)a_{\alpha_t}(u,u)\leq A_t(u)
 \leq(1+Ct)a_{\alpha_t}(u,u).
\end{equation}
Also
\begin{equation}\label{eq:L2-compare}
 \int\|U_zu\|_{\Ltwo}^2\eta_t(z)\,dz
 \leq(1+Ct)\|u\|_{\Ltwo}^2.
\end{equation}

Average \eqref{eq:ehrling} with \(v=U_zu\) and use
\eqref{eq:mass-average}--\eqref{eq:L2-compare}.  Absorbing the harmless
volume component in \eqref{eq:positive-mass}, we obtain the central estimate:
for every \(\varepsilon>0\), there is \(C_\varepsilon'\), independent of all
sufficiently small \(t\), such that
\begin{equation}\label{eq:uniform-ehrling}
 b_{\mu_t}(u,u)\leq
 \varepsilon a_{\alpha_t}(u,u)
 +C_\varepsilon'\|u\|_{\Ltwo}^2
 \qquad(u\in H^1(M)).
\end{equation}
The estimate initially holds for smooth \(u\) and extends by density because
\(\alpha_t\) and the density of \(\mu_t\) are smooth and positive.

For \(p=2\) and \(\alpha=1\), retain \(a_1\) instead of defining
\(a_{\alpha_t}\).  Equation \eqref{eq:metric-close} compares
\(\int a_1(U_zu)\eta_t dz\) directly with \(a_1(u)\), so the same estimate
\eqref{eq:uniform-ehrling} holds with \(\alpha_t=1\).

\subsection*{3. Convergence of the energy and mass forms}

Take any sequence \(t_n\downarrow0\), and abbreviate
\(a_n=a_{\alpha_{t_n}}\), \(b_n=b_{\mu_{t_n}}\).  By
\eqref{eq:Lr}, \(\alpha_{t_n}\to\alpha\) in \(L^1\).  The energy forms have
the following two elementary Mosco properties:
\begin{enumerate}[label=(\alph*)]
\item If \(u_n\to u\) in \(\Ltwo\) and
      \(\sup_n a_n(u_n,u_n)<\infty\), then
\begin{equation}\label{eq:liminf-energy}
 a_\alpha(u,u)\leq\liminf_n a_n(u_n,u_n).
\end{equation}
\item For every \(u\in H_\alpha\), there are smooth \(u_n\to u\) in
      \(\Ltwo\) such that
\begin{equation}\label{eq:recovery-energy}
 a_n(u_n,u_n)\longrightarrow a_\alpha(u,u).
\end{equation}
\end{enumerate}
For (a), \eqref{eq:lower} bounds \(u_n\) in ordinary \(H^1\).  After taking
a subsequence, \(du_n\rightharpoonup du\) in \(L^2\), while
\(\sqrt{\alpha_{t_n}}\to\sqrt\alpha\) in \(L^2\).  The bounded fields
\(\sqrt{\alpha_{t_n}}du_n\) therefore converge weakly to
\(\sqrt\alpha,du\); weak lower semicontinuity proves
\eqref{eq:liminf-energy}.  For (b), approximate \(u\) in \(H_\alpha\) by
smooth functions and choose a diagonal sequence; for each fixed smooth
function, \(L^1\)-convergence of the coefficients gives energy convergence.
The same diagonal argument works simultaneously on any fixed
finite-dimensional subspace and preserves all polarized energy products.

The uniform estimate upgrades this to convergence of the moving mass forms.
Precisely, if
\begin{equation}\label{eq:moving-assumption}
 u_n\to u\text{ in }\Ltwo,\qquad \sup_n a_n(u_n,u_n)<\infty,
\end{equation}
then
\begin{equation}\label{eq:moving-mass}
 b_n(u_n,u_n)\longrightarrow b_\mu(u,u).
\end{equation}
To see this, choose smooth \(v_j\to u\) in \(H_\alpha\), with their energies
uniformly bounded.  For fixed \(j\), weak-star convergence gives
\(b_n(v_j,v_j)\to b_\mu(v_j,v_j)\).  On the other hand,
\eqref{eq:uniform-ehrling} and
\(a_n(u_n-v_j)\leq2a_n(u_n)+2a_n(v_j)\) give, for every fixed
\(\varepsilon>0\),
\[
 \limsup_n b_n(u_n-v_j,u_n-v_j)
 \leq C\varepsilon+C_\varepsilon'\|u-v_j\|_{\Ltwo}^2,
\]
where \(C\) is independent of \(j\) for large \(j\).  First choose
\(\varepsilon\) small and then \(j\) large.  The triangle inequality in
\(L^2(\mu_n)\), followed by the corresponding limiting argument in
\(L^2(\mu)\), proves \eqref{eq:moving-mass}.  Polarization gives convergence
of mixed mass products for finitely many such sequences.

We also need a uniform coercivity estimate.  There is \(C<\infty\) such that
\begin{equation}\label{eq:coercivity}
 \|u\|_{\Ltwo}^2\leq C\bigl(a_n(u,u)+b_n(u,u)\bigr)
\end{equation}
for all large \(n\).  Indeed, subtract the volume mean \(\bar u\).  The
ordinary Poincar\'e inequality and \eqref{eq:lower} bound
\(\|u-\bar u\|_{\Ltwo}^2\) by \(Ca_n(u,u)\).  Estimate
\eqref{eq:uniform-ehrling} then bounds \(b_n(u-\bar u,u-\bar u)\) by the same
quantity.  Finally,
\[
 \mu(M)^{1/2}|\bar u|
 =\|\bar u\|_{L^2(\mu_n)}
 \leq\|u\|_{L^2(\mu_n)}+
       \|u-\bar u\|_{L^2(\mu_n)},
\]
which proves \eqref{eq:coercivity}.

\subsection*{4. Two-sided min--max convergence}

For the limsup, let \(E\subset H_\alpha\) be the span of the first
\(k+1\) limiting eigenfunctions.  The finite-dimensional version of the
recovery construction supplies maps \(R_n:E\to C^\infty(M)\) for which all
energy products converge.  Equations \eqref{eq:moving-mass} and polarization
give convergence of all mass products as well.  The two Gram matrices thus
converge uniformly on the unit sphere of \(E\), and the min--max
principle~\cite{Davies} gives
\begin{equation}\label{eq:limsup}
 \limsup_n\lambda_k(\alpha_{t_n},\mu_{t_n})
 \leq\lambda_k(\alpha,\mu).
\end{equation}

For the reverse inequality, pass to a subsequence realizing the liminf.  The
claim is immediate if the liminf is infinite, so assume it is finite.  Let
\(u_{n,0},\ldots,u_{n,k}\) be first eigenfunctions for the smooth pair,
chosen \(b_n\)-orthonormal.  Their energies are at most
\(\lambda_k(\alpha_{t_n},\mu_{t_n})\).  Estimate
\eqref{eq:coercivity} bounds them in \(L^2(v_g)\), and
\eqref{eq:lower} bounds them in ordinary \(H^1\).  Rellich compactness and a
diagonal extraction give
\[
 u_{n,j}\longrightarrow u_j\quad\text{in }L^2(v_g),
 \qquad j=0,\ldots,k.
\]
Equation \eqref{eq:moving-mass} and polarization show that the limits are
\(b_\mu\)-orthonormal.  For every \(c=(c_0,\ldots,c_k)\), energy lower
semicontinuity gives
\begin{align*}
 a_\alpha\left(\sum_jc_ju_j,\sum_jc_ju_j\right)
 &\leq\liminf_n
 a_n\left(\sum_jc_ju_{n,j},\sum_jc_ju_{n,j}\right)\\
 &\leq\left(\liminf_n\lambda_k(\alpha_{t_n},\mu_{t_n})\right)
       \sum_j|c_j|^2.
\end{align*}
The span of the \(u_j\)'s is therefore an admissible \((k+1)\)-dimensional
space in the limiting min--max problem, and
\begin{equation}\label{eq:liminf}
 \lambda_k(\alpha,\mu)
 \leq\liminf_n\lambda_k(\alpha_{t_n},\mu_{t_n}).
\end{equation}
Together, \eqref{eq:limsup} and \eqref{eq:liminf} prove eigenvalue
convergence.  Property \eqref{eq:Lr}, preservation of \(\mu(M)\), and the
special \(p=2\) construction prove the remaining assertions of
Theorem~\ref{thm:main}. \qed

\section*{Verification notes}

The proof was checked against four common failure modes.
\begin{enumerate}[leftmargin=*]
\item Weak-star convergence alone gives only the wrong min--max direction.
      The proof never relies on it alone; \eqref{eq:uniform-ehrling} supplies
      collective compactness.
\item Translating only \(\mu\) is not compatible with a rough unbounded
      \(\alpha\).  Here both measures are averaged by the same maps, and
      \eqref{eq:energy-compare} is a multiplicative comparison independent of
      pointwise upper bounds on \(\alpha\).
\item Strong \(L^1\) convergence of coefficients does not by itself imply
      operator-norm form convergence.  The proof uses the weaker but adequate
      Mosco liminf plus finite-dimensional recovery properties.
\item Averaging spectral inequalities can turn a rank-\(k\) exceptional space
      into an infinite-rank one.  The argument avoids averaging spectral
      projectors: it averages the compactness inequality and applies min--max
      only after compact convergence has been proved.
\end{enumerate}

The most important points for human review are the smooth-kernel consequence
of the spanning-flow construction, the uniform estimate
\eqref{eq:uniform-ehrling}, and the diagonal finite-dimensional recovery used
in \eqref{eq:limsup}.

\section*{Novelty check and limitations}

A bounded search was performed before promotion.  Local run indexes were
searched for arXiv:\allowbreak 2601.17896, the exact Remark~1.3 wording,
compact weighted
traces, smooth-measure approximation, Mosco convergence, and spectral
convergence.  The source of Girouard--Karpukhin--Lagac\'e arXiv:2004.10784 was
also searched.  A web pass used the exact question phrase and combinations of
``joint mollification'', ``weighted Dirichlet form'', ``compact trace'', and
``eigenvalue approximation''.  It found general Mosco-convergence literature
but no theorem matching Theorem~\ref{thm:main}, and no later paper explicitly
answering Vinokurov's question.  Novelty confidence is therefore moderate,
not definitive.

The theorem does not include weights that can vanish.  If Remark~1.3 is read
as asking simultaneously about arbitrary degenerate \(\alpha\), this packet
should be reclassified as a sharp partial result.  Under the uniformly
positive hypothesis stated immediately before the question, it is a full
approximation theorem.

\section*{Human-review recommendation}

Send to an analyst working on weighted Dirichlet forms or spectral geometry.
The argument is short enough to verify directly and appears to close the
uniformly positive case, but the interpretation of the source's one-argument
notation \(H^1(M,\alpha)\) should be confirmed with the author before a
publication claim.

\begin{thebibliography}{9}

\bibitem{Vinokurov}
D. Vinokurov,
\emph{Eigenvalue optimization in higher dimensions and \(p\)-harmonic maps},
arXiv:2601.17896v3 (2026), especially Remark~1.3 and
Section~2.6.

\bibitem{Lee}
J. M. Lee,
\emph{Introduction to Smooth Manifolds}, 2nd ed.,
Graduate Texts in Mathematics 218, Springer, 2013.

\bibitem{Davies}
E. B. Davies,
\emph{Spectral Theory and Differential Operators},
Cambridge Studies in Advanced Mathematics 42, Cambridge University Press,
1995.

\end{thebibliography}

\end{document}
